\documentclass[11pt,a4paper]{article}
\usepackage[margin=25mm]{geometry}
\usepackage{amsmath,amssymb,amsthm,booktabs}
\newtheorem{lemma}{Lemma}\newtheorem{theorem}{Theorem}
\title{Chiral quotients with small regular covers}
\author{Human Compression project}\date{}
\begin{document}\maketitle
Hartley, Hubard and Leemans, \emph{Two atlases of abstract chiral
polytopes for small groups}, Ars Mathematica Contemporanea 5 (2012), 371--382, report
56 chiral polytopes whose minimal regular covers have at most 2000 flags,
excluding covers with 1024 or 1536 flags. We verify the 56 listed examples
using the literal companion data acquired on 15 October 2026. This note
proves their existence and distinctness. Exhaustiveness among all covers
is a separate obligation and is not asserted here. A distinct dual is
counted separately, while the two enantiomorphic markings describe one
unrooted polytope.

For each entry the retained data specify permutations
$r_0,\ldots,r_{d-1}$ generating $W$, and words generating a subgroup
$N\leq W$. Put $H=\langle r_0r_i:1\leq i<d\rangle$.
The vertices of the flag graph are the right cosets $Nw$; the edge of
colour $i$ sends $Nw$ to $Nwr_i$. These operations are unambiguous.
We acquire the previously proved normalizer theorem: the group of
automorphisms preserving every colour is $N_W(N)/N$. In particular,
if $N_W(N)=H$ has index two and avoids every $r_i$, the graph has two
automorphism orbits, and each adjacency exchanges them. To obtain a
polytope one must also verify the incidence axioms; normalizer data
alone do not imply them.

\begin{lemma}[Reconstruction from a finite flag graph]
Let a finite graph have $d$ specified involutory adjacencies, none of
which fixes a vertex. For each
$i$, call the connected components using all colours except $i$ its
faces of rank $i$. Associate to each graph vertex its tuple of faces.
Suppose the following finite checks hold: these tuples are distinct;
incidence defined by occurrence in a tuple is transitive; its maximal
chains are exactly the tuples; every section of rank two has exactly
two intermediate faces; and the graph vertices containing any fixed
chain are connected by the remaining colours. Then adjoining a least
and a greatest face gives an abstract polytope with precisely the
specified flag graph.
\end{lemma}
\begin{proof}
Ranks orient incidence between distinct ranks. Transitivity, together
with equality within a rank, gives a partial order. Each face occurs
in a tuple, and each comparable pair does so by definition. Every chain
extends in the finite order to a maximal chain, which by hypothesis is
a full tuple. Thus all maximal chains have the required length.
The section checks are exactly the diamond axiom, including sections
incident with the least or greatest face. The connectedness hypothesis
is strong flag connectivity. An edge of colour $i$ preserves every
other component, and its endpoints differ in the remaining face by
injectivity of tuples. The diamond axiom makes this the unique adjacent
flag of that colour. Hence the given graph is exactly the flag graph.
\end{proof}

\begin{lemma}[Complete comparison of coloured graphs]
Starting at a chosen vertex, traverse a connected graph in breadth-first
order, considering colours in a fixed order and numbering each vertex
when it is first encountered. The resulting numbered adjacency table
is a complete invariant of the graph with its chosen root. Trying each
possible image of a root therefore finds all colour-preserving
isomorphisms. The least rooted table is a complete unrooted invariant.
\end{lemma}
\begin{proof}
A rooted colour-preserving isomorphism must carry the first traversal
to the second, step by step, so it preserves all assigned numbers and
the resulting table. Conversely, equality of tables gives the explicit
bijection between equally numbered vertices, preserving every edge.
Connectedness makes an isomorphism unique once its root image is fixed.
Taking every possible root proves both assertions.
\end{proof}

\begin{theorem}[The listed examples]
The retained 56 entries give 48 pairwise nonisomorphic chiral polytopes
of rank three and eight of rank four. Each has its stated face numbers,
regular monodromy cover and rotation-group SmallGroup identifier.
\end{theorem}
\begin{proof}
The literal permutation groups have the source orders and SmallGroup
identifiers. For every pair of subsets of generators the checker verifies
the complete intersection property. It also checks involutions, string
commutation and each adjacent product order. Thus these are the stated
regular polytopes, using the retained string C-group correspondence.

The stated words generate subgroups of exactly the stated orders. In
each case $N\leq H$, $[W:H]=2$, $N_W(N)=H$, and
$\operatorname{Core}_W(N)=1$. The quotient $H/N$ has the stated
SmallGroup identifier, verified on that actual quotient. The core test
makes the flag action faithful, so its monodromy group is the literal
group $W$. The normalizer theorem identifies the full automorphism
group and its two flag orbits.

For all 19,712 cosets the checker constructs the complete adjacency
table. It computes every face component, verifies incidence transitivity,
enumerates every maximal chain and checks all diamonds. For each subset
of ranks it groups flags by their common faces and compares the resulting
sets with components in the complementary colours. Thus every chain
residue is connected, and the reconstruction lemma applies. The computed
face counts agree with every source entry.

Finally the second lemma checks every possible root image directly from
the graphs, without using the group normalizer calculation. For each
graph the images of a fixed root under its automorphisms are exactly
the even vertices in its bipartition. Their number equals the stated
rotation-group order. The 56 least adjacency tables are distinct.
This proves chirality, full automorphism orders and the stated
equivalence convention independently at the graph level.
\end{proof}

\paragraph{The rank-four rows.}
The four entries below and their reversals are the eight rank-four
examples. We write $(n,j)$ for SmallGroup$(n,j)$; a group order alone
is never used to identify a group.
\begin{center}\small
\begin{tabular}{cccr}\toprule
Type&Regular cover group&Rotation group&$|N|$\\\midrule
$\{3,3,8\}$&$(768,1086052)$&$(192,1494)$&2\\
$\{3,6,9\}$&$(972,108)$&$(162,19)$&3\\
$\{3,6,18\}$&$(1944,949)$&$(324,77)$&3\\
$\{6,6,9\}$&$(1944,949)$&$(324,77)$&3\\\bottomrule
\end{tabular}\end{center}

\paragraph{Certificates, attribution and remaining scope.}
The pages are parsed as literal integers, disjoint cycles and group
words; no downloaded program is executed. Exact source bytes, generated
calls and terminal receipts are retained. GAP's finite permutation-group,
normalizer, coset-action and SmallGroup identification algorithms remain
explicit software premises. The incidence and isomorphism checks use
only the complete finite adjacency tables and ordinary integer operations.
The rotation relators displayed on the pages are also satisfied, but
this does not prove that every displayed abstract presentation has no
additional elements. Presentation equality is not claimed by that test.

The source authors retain credit for the examples and classification.
The normalizer argument, incidence reconstruction and deterministic
rooted comparison supply a uniform verification. No new open-question
solution or compression measurement is asserted. The earlier integral
proof for the order-588 example is acquired separately. Present website
bytes are not asserted to be identical to the 2012 website. A complete
reconstruction of the source still needs exclusion of all other regular
covers, and its separate almost-simple atlas.
\end{document}
